\chapter{Introduction to Python, Variables, Data Types \& Operators}

Welcome to Programming in Python! Python is the undisputed language of modern Data Science, Machine Learning, and Artificial Intelligence. Its clean, human-readable syntax makes it easy to write code, while its powerful ecosystem of libraries (like NumPy, Pandas, and PyTorch) allows data scientists to build cutting-edge intelligent systems.

\section{Variables and Literals}

In Python, a \textbf{variable} is a named location in memory used to store data. Unlike languages like C++ or Java, Python uses \textbf{dynamic typing}, meaning you do not need to explicitly declare a variable's data type when creating it.

\begin{definitionbox}{Variables and Assignment}
A variable is created the moment you first assign a value to it using the assignment operator (`=`).
\begin{lstlisting}
x = 10          # An integer variable
name = "IITM"   # A string variable
pi = 3.14159    # A floating-point variable
is_valid = True # A boolean variable
\end{lstlisting}
\end{definitionbox}

\textbf{Data Science Relevance:} In data science, variables hold feature columns, target labels, model hyperparameters (like learning rate $\alpha = 0.01$), and dataset paths.

\begin{videocard}{IITM BS Lecture: Introduction to Variables}{https://www.youtube.com/watch?v=Yg6xzi2ie5s}
Official lecture introducing Python variables, assignment statements, and memory concepts.
\end{videocard}

\begin{videocard}{IITM BS Lecture: Variables and Literals}{https://www.youtube.com/watch?v=tDaXdoKfX0k}
Official lecture explaining literals, variable naming rules, and assignment mechanics.
\end{videocard}

\section{Core Primitive Data Types}

Python has four primary built-in primitive data types:

\begin{definitionbox}{Primitive Data Types}
\begin{itemize}
    \item \textbf{Integer (`int`):} Whole numbers without decimals (e.g., `-5`, `0`, `42`).
    \item \textbf{Float (`float`):} Real numbers with decimal points or scientific notation (e.g., `3.14`, `-0.001`, `2.5e-4`).
    \item \textbf{String (`str`):} Sequences of text characters enclosed in single, double, or triple quotes (e.g., `"data"`, `'Python'`).
    \item \textbf{Boolean (`bool`):} Truth values: `True` or `False`.
\end{itemize}
You can inspect the type of any variable using the built-in `type()` function.
\end{definitionbox}

\begin{examplebox}{Checking Types and Type Casting}
\begin{lstlisting}
age = 25
print(type(age))  # Output: <class 'int'>

# Type Conversion (Type Casting)
height_str = "1.75"
height_float = float(height_str)  # Converts string "1.75" to float 1.75
print(type(height_float))         # Output: <class 'float'>
\end{lstlisting}
\end{examplebox}

\begin{videocard}{IITM BS Lecture: Data Types (Parts 1 \& 2)}{https://www.youtube.com/watch?v=8n4MBjuDBu4}
Official lectures covering numeric data types, type conversion, and character strings.
\end{videocard}

\section{String Slicing and Basic String Operations}

Strings are immutable sequences of characters. You can access individual characters using 0-based indexing or slice sub-strings using the bracket syntax `[start:stop:step]`.

\begin{formulabox}{String Slicing Syntax}
Given `s = "DataScience"`:
\begin{itemize}
    \item `s[0]` $\to$ `'D'` (First character)
    \item `s[-1]` $\to$ `'e'` (Last character via negative indexing)
    \item `s[0:4]` $\to$ `'Data'` (Slice from index 0 up to, but not including, index 4)
    \item `s[::-1]` $\to$ `'ecneicSataD'` (Reverse the string)
\end{itemize}
\end{formulabox}

\begin{videocard}{IITM BS Lecture: Strings and String Slicing}{https://www.youtube.com/watch?v=sS89tiDuqoM}
Official lecture demonstrating string indexing, negative indices, and slicing notation.
\end{videocard}

\section{Operators and Expressions}

Python supports arithmetic, comparison (relational), logical, and assignment operators.

\begin{theorembox}{Arithmetic and Logical Operators}
\begin{itemize}
    \item \textbf{Arithmetic:} `+` (add), `-` (subtract), `*` (multiply), `/` (float division), `//` (floor division), `%` (modulus / remainder), `\textbf{` (exponentiation).
    \item \textbf{Comparison:} `==` (equal), `!=` (not equal), `>`, `<`, `>=`, `<=`.
    \item \textbf{Logical:} `and`, `or`, `not`.
\end{itemize}
\end{theorembox}

\begin{examplebox}{Floor Division vs Division}
\begin{lstlisting}
a = 17
b = 5

print(a / b)   # 3.4   (Float Division)
print(a // b)  # 3     (Floor Division: drops decimal part)
print(a % b)   # 2     (Modulus: remainder of 17 / 5)
print(a } b)  # 1419857 (17 raised to the power of 5)
\end{lstlisting}
\end{examplebox}

\begin{videocard}{IITM BS Lecture: Operators and Expressions}{https://www.youtube.com/watch?v=8pu73HKzNOE}
Official lecture detailing operator precedence, modulus, floor division, and logical expressions.
\end{videocard}

\newpage
\section{Practice Exercises}
\textit{Detailed step-by-step solutions are available in the Appendix.}

\begin{enumerate}
    \item \textbf{Variable Types:} Identify the data type (`int`, `float`, `str`, or `bool`) of the following Python expressions:
    \begin{enumerate}
        \item `x = 10 / 2`
        \item `y = 10 // 2`
        \item `z = "100" + "200"`
        \item `w = (5 > 3) and (2 == 4)`
    \end{enumerate}

    \item \textbf{Arithmetic Logic:} Calculate the exact numerical output of the following Python expressions without running the code:
    \begin{enumerate}
        \item `2 + 3 * 4 ** 2`
        \item `25 % 7`
        \item `19 // 4`
    \end{enumerate}

    \item \textbf{String Slicing:} Given the string `text = "MachineLearning"`:
    \begin{enumerate}
        \item Write a slice to extract `"Machine"`.
        \item Write a slice to extract `"Learning"`.
        \item What is the output of `text[-8:]`?
        \item Write a slice to reverse the entire string.
    \end{enumerate}

    \item \textbf{Type Casting & Input:} A user enters their age as a string `"22"`. Write a 2-line Python script that converts this input into an integer, adds $5$ years, and prints: `"In 5 years, you will be 27"`.

    \item \textbf{Data Science Application:} You are calculating the Body Mass Index (BMI) of a patient. The formula is:
    $$\text{BMI} = \frac{\text{weight in kg}}{(\text{height in meters})^2}$$
    Given `weight = 75` and `height = 1.75`, write a Python snippet to compute the BMI and round it to 2 decimal places using `round(bmi, 2)`.
\end{enumerate}